\documentclass[withoutpreface,bwprint]{cumcmthesis} %去掉封面与编号页，电子版提交的时候使用。


\usepackage[framemethod=TikZ]{mdframed}
\usepackage{url}   % 网页链接
\usepackage{subcaption} % 子标题
\title{全国大学生数学建模竞赛编写的 \LaTeX{} 模板}
\tihao{A}
\baominghao{4321}
\schoolname{XX大学}
\membera{ }
\memberb{ }
\memberc{ }
\supervisor{ }%辅导老师
\yearinput{2023}
\monthinput{9}
\dayinput{8}

\begin{document}

\maketitle
\begin{abstract}

开头段：需要充分概括论文内容，一般两到三句话即可，长度控制在三至五行。

问题一中，解决了什么问题；应用了什么方法；得到了什么结果。

问题二中，解决了什么问题；应用了什么方法；得到了什么结果。

问题三中，解决了什么问题；应用了什么方法；得到了什么结果。

结尾段：可以总结下全文，也可以介绍下你的论文的亮点，也可以对类似的问题进行适当的推广。

\begin{mdframed} [%
	roundcorner=5pt,
	linecolor=gray!50,
	outerlinewidth=0.5pt,
	middlelinewidth=0.3pt, backgroundcolor=gray!2,
innertopmargin=\topskip, frametitle={首页三要素: 论文标题 + 摘要 + 关键词},
frametitlefont= \bfseries,frametitlerule=true,frametitlealignment =\raggedright\noindent,
frametitlerulewidth=.5pt, frametitlebackgroundcolor=gray!2,]

标题：

基于所使用的主要模型或者方法作为标题（推荐）

直接使用赛题所给的题目或者要研究的问题作为标题

摘要：

摘要是数模论文写作中最重要的一部分，因为评阅老师的时间有限，拿到一篇论文后不会完整的从头读到尾，所以评阅老师往往会重点阅读摘要部分，并结合官方的评阅要点来对你的论文进行初步评定。因此，大家一定要好好打磨论文的摘要，摘要一般是其他部分都完成后再来书写，写完后需要反复阅读反复修改。

关键词：

关键词一般放4-6个，可以放论文中使用的主要模型，也可以放论文里面出现次数较多，能体现论文的主要内容的词。

\end{mdframed}

\keywords{关键词1\quad  关键词2\quad   关键词3\quad  关键词4}
\end{abstract}

\section{问题重述}

数学建模比赛论文是要我们解决一道给定的问题，所以正文部分一般应从问题重述开始，一般确定选题后就可以开始写这一部分了。

这部分的内容是将原问题进行整理，将问题背景和题目分开陈述即可，所以基本没啥难度。

本部分的目的是要吸引读者读下去，所以文字不可冗长，内容选择不要过于分散、琐碎，措辞要精练。

注意：在写这部分的内容时，绝对不可照抄原题！（论文会查重）

应为：在仔细理解了问题的基础上，用自己的语言重新将问题描述一遍。语言需要简明扼要，没有必要像原题一样面面俱到。

\section{问题分析}

\subsection{问题一的分析}

从实际问题到模型建立是一种从具体到抽象的思维过程，问题分析这一部分就是沟通这一过程的桥梁，因为它反映了建模者对于问题的认识程度如何，也体现了解决问题的雏形，起着承上启下的作用，也很能反应出建模者的综合水平。

这部分的内容应包括：题目中包含的信息和条件，利用信息和条件对题目做整体分析，确定用什么方法建立模型，一般是每个问题单独分析一小节，分析过程要简明扼要，不需要放结论。

建议在文字说明的同时用图形或图表（例如流程图）列出思维过程，这会使你的思维显得很清晰，让人觉得一目了然。

\subsection{问题二的分析}

\subsection{问题三的分析}

\section{模型假设}

\begin{enumerate}
    \item 题目明确给出的假设条件
    \item 排除生活中的小概率事件(例如黑天鹅事件、非正常情况)
    \item 仅考虑问题中的核心因素，不考虑次要因素的影响
    \item 使用的模型中要求的假设
    \item 对模型中的参数形式(或者分布)进行假设
    \item 和题目联系很紧密的一些假设，主要是为了简化模型
\end{enumerate}

\section{符号说明}

\begin{table}[!htbp]
    \caption{标准三线表格}\label{tab:001} \centering
    \begin{tabular}{ccccc}
        \toprule[1.5pt]
        符号 & 说明 & 单位\\
        \midrule[1pt]
        5 & 269.8 & 0.000674\\
        10 & 421.0 & 0.001035\\
        20 & 640.2 & 0.001565\\
        \bottomrule[1.5pt]
    \end{tabular}
\end{table}

本部分是对模型中使用的重要变量进行说明，一般排版时要放到一张表格中。

注意：第一：不需要把所有变量都放到这个表里面，模型中用到的临时变量可以不放。第二：下文中首次出现这些变量时也要进行解释，不然会降低文章的可读性。

\section{模型的建立与求解}

\subsection{问题一模型的建立与求解}

\subsubsection{模型的建立}

模型建立是将原问题抽象成用数学语言的表达式，它一定是在先前的问题分析和模型假设的基础上得来的。因为比赛时间很紧，大多时候我们都是使用别人已经建立好的模型。这部分一定要将题目问的问题和模型紧密结合起来，切忌随意套用模型。我们还可以对已有模型的某一方面进行改进或者优化，或者建立不同的模型解决同一个问题，这样就是论文的创新和亮点。

\subsubsection{模型的求解}

把实际问题归结为一定的数学模型后，就要利用数学模型求解所提出的实际问题了。一般需要借助计算机软件进行求解，例如常用的软件有Matlab, Spss, Lingo, Excel, Stata, Python等。求解完成后，得到的求解结果应该规范准确并且醒目，若求解结果过长，最好编入附录里。（注意：如果使用智能优化算法或者数值计算方法求解的话，需要简要阐明算法的计算步骤）

\subsection{问题二模型的建立与求解}

\begin{equation}
E=mc^2
\label{eq:energy}
\end{equation}
式\cref{eq:energy}是质能方程。

\subsection{问题三模型的建立与求解}

\begin{theorem}
    这是一个定理。
    \label{thm:example}
\end{theorem}
由\cref{thm:example}我们知道了定理环境的使用。

\section{模型的分析与检验}

模型的分析与检验的内容也可以放到模型的建立与求解部分，这里我们单独抽出来进行讲解，因为这部分往往是论文的加分项，很多优秀论文也会单独抽出一节来对这个内容进行讨论。

模型的分析 ：在建模比赛中模型分析主要有两种，一个是灵敏度(性)分析，另一个是误差分析。灵敏度分析是研究与分析一个系统（或模型）的状态或输出变化对系统参数或周围条件变化的敏感程度的方法。其通用的步骤是：控制其他参数不变的情况下，改变模型中某个重要参数的值，然后观察模型的结果的变化情况。误差分析是指分析模型中的误差来源，或者估算模型中存在的误差，一般用于预测问题或者数值计算类问题。

模型的检验：模型检验可以分为两种，一种是使用模型之前应该进行的检验，例如层次分析法中一致性检验，灰色预测中的准指数规律的检验，这部分内容应该放在模型的建立部分；另一种是使用了模型后对模型的结果进行检验，数模中最常见的是稳定性检验，实际上这里的稳定性检验和前面的灵敏度分析非常类似，等会大家看到例子就明白了。(大家尽量在论文中使用灵敏度分析，视频中有详细的讲解)

\section{模型的评价与推广}

注：本部分的标题需要根据你的内容进行调整，例如：如果你没有写模型推广的话，就直接把标题写成模型的评价与改进。很多论文也把这部分的内容直接统称为“模型评价”部分，也是可以的。

\subsection{模型的优点}

优缺点是必须要写的内容，改进和推广是可选的，但还是建议大家写，实力比较强的建模者可以在这一块充分发挥，这部分对于整个论文的作用在于画龙点睛。

\subsection{模型的缺点}

缺点写的个数要比优点少

\subsection{模型的改进}

主要是针对模型中缺点有哪些可以改进的地方

\subsection{模型的推广}

将原题的要求进行扩展，进一步讨论模型的实用性和可行性

\section{参考文献}

所有引用他人或公开资料(包括网上资料)的成果必须按照科技论文的规范列出参考文献，并在正文引用处予以标注。

常见的三种参考文献的表达方式（标准不唯一）：

书籍的表述方式为： [编号] 作者，书名，出版地：出版社，出版年月。

期刊杂志论文的表述方式为： [编号] 作者，论文名，杂志名，卷期号：起止页码，出版年。

网上资源(例如数据库、政府报告)的表述方式为： [编号] 作者，资源标题，网址，访问时间。

%参考文献
\begin{thebibliography}{9}%宽度9
    \bibitem[1]{liuhaiyang2013latex}
    刘海洋.
    \newblock \LaTeX {}入门\allowbreak[J].
    \newblock 电子工业出版社, 北京, 2013.
    \bibitem[2]{mathematical-modeling}
    全国大学生数学建模竞赛论文格式规范 (2023 年 修改).
    \bibitem{3} \url{https://www.latexstudio.net}
\end{thebibliography}

\newpage
%附录
\begin{appendices}

除了支撑材料的文件列表和源程序代码外，附录中还可以包括下面内容：

\begin{itemize}
    \item 某一问题的详细证明或求解过程；
    \item 自己在网上找到的数据；
    \item 比较大的流程图；
    \item 较繁杂的图表或计算结果
\end{itemize}

\section{介绍：支撑材料的文件列表}

\section{介绍：该代码是某某语言编写的，作用是什么}

\begin{lstlisting}[language=python]
import sympy as sp
import numpy as np
import matplotlib.pyplot as plt


t = sp.Symbol('t')
r = sp.Symbol('r')
x0 = sp.Symbol('x0')
x = sp.Function('x')(t)

# 求解微分方程 Dx = r*x, x(0) = x0
eq = sp.Eq(x.diff(t), r * x)
solution = sp.dsolve(eq, ics={x.subs(t, 0): x0})
print("马尔萨斯模型通解:", solution)

# 替换x0和r为确定值
x0_val = 1000
r_val = 0.1
x_specific = solution.subs({x0: x0_val, r: r_val})
print(f"当x0={x0_val}, r={r_val}时的解:", x_specific)

t_range = np.linspace(0, 50, 1000)  # 0到50年，取1000个点

# 马尔萨斯模型公式：x(t) = x0 * exp(r*t)
x_vals = x0_val * np.exp(r_val * t_range)
plt.plot(t_range, x_vals, label=f'growth rate = {r_val*100:.1f}%')
plt.title('Malthus Population Growth Model')
plt.xlabel('Time (Years)')
plt.ylabel('Population Size')
plt.grid(True)
plt.legend()
plt.tight_layout()
plt.show()


# 定义符号
xm = sp.Symbol('xm')  # 环境容纳量
t0 = sp.Symbol('t0')  # 初始时间

# 求解微分方程 Dx = r*(1 - x/xm)*x, x(t0) = x0
eq_logistic = sp.Eq(x.diff(t), r * (1 - x/xm) * x)
solution_logistic = sp.dsolve(eq_logistic, ics={x.subs(t, t0): x0})
print("Logistic模型通解:", solution_logistic)

# 代入具体参数
t0_val = 0
x0_log = 1000
xm_val = 10000
r_log = 0.05
logistic_specific = solution_logistic.subs({t0: t0_val, x0: x0_log, xm: xm_val, r: r_log})
print("代入参数后的Logistic模型解:", logistic_specific)

# 将符号表达式转换为可计算的函数
logistic_func = sp.lambdify(t, logistic_specific.rhs, 'numpy')
t_logistic = np.linspace(0, 200, 1000)
x_logistic = logistic_func(t_logistic)

plt.figure(figsize=(10, 6))
plt.plot(t_logistic, x_logistic, '-', label='Logistic Growth', color='orange')
plt.title('Logistic Growth Model')
plt.xlabel('Time (Years)')
plt.ylabel('Population Size')
plt.grid(True)
plt.legend()
plt.tight_layout()
plt.show()
\end{lstlisting}

\section{介绍：该代码是某某语言编写的，作用是什么}

\begin{lstlisting}[language=python]
import numpy as np
from scipy.optimize import linprog

# 1. 定义目标函数（最大化价值，转为最小化负价值）
c = np.array([-540, -200, -180, -350, -60, -150, -280, -450, -320, -120])

# 2. 定义约束条件（重量不超过30）
# 不等式约束：A·x ≤ b（物品重量约束）
A = np.array([[6, 3, 4, 5, 1, 2, 3, 5, 4, 2]])  # 物品重量矩阵
b = np.array([30])  # 最大承重

# 3. 变量范围（0-1整数变量）
# 注意：scipy的linprog不直接支持整数约束，需后续处理为0-1
lb = [0] * 10  # 下限：0
ub = [1] * 10  # 上限：1

# 4. 求解线性规划（松弛问题，先不强制整数）
# 由于scipy无专门整数规划函数，先求解连续解，再四舍五入为0-1
result = linprog(c, A, b, bounds=list(zip(lb, ub)), method='highs')

# 5. 处理结果（将连续解转为0-1整数解）
# 四舍五入（接近1的取1，接近0的取0）
x_int = np.round(result.x).astype(int)

# 6. 计算实际目标函数值（最大化的总价值）
# 原始价值向量
values = np.array([540, 200, 180, 350, 60, 150, 280, 450, 320, 120])
total_value = np.sum(x_int * values)

# 7. 输出结果
print("决策变量（0-1）：", x_int)
print("选中的物品索引：", np.where(x_int == 1)[0] + 1)  # 转为1-based索引
print("最大总价值：", total_value)
print("总重量：", np.sum(x_int * A[0]))  # 验证是否满足重量约束
\end{lstlisting}
\end{appendices}

\end{document} 